% GENERATED FILE. Edit canonical lessons, then run npm run generate:books.
% canonical-source-hash: fnv1a64:1ca0e82c08f24c5d
% canonical-lessons: FR-C192-cuit-cuite, FR-C192-mur-mure, FR-C192-sombre, FR-C192-clair-claire, FR-C192-lointain-lointaine, FR-R192-first-pass-words-from-chapters-184-187, FR-R192-second-pass-words-from-chapters-188-192

\chapter{Cooked, Ripe, Dark, Clear, Distant}
\label{ch:fr-cooked-ripe-dark-clear-distant}
\hlchaptermodality{192}

\begin{chapteropening}
\textbf{By the end of this chapter:} \emph{I can say cooked, ripe, dark, clear and distant.}

Complete the last lesson of chapter 192: I can say cooked, ripe, dark, clear and distant.
\end{chapteropening}

\section[cuit / cuite]{cuit / cuite — cooked}

\label{lesson:FR-C192-cuit-cuite}



\begin{quote}
Before the new one: say the French for useful, then the French for useless.
\end{quote}

\subsection*{You'll want to know: cuit / cuite}
\textbf{cuit / cuite} — \textquotedblleft{}cooked\textquotedblright{}.

\begin{grammarlens}[title={What you've built}]
One more word for this chapter.
\end{grammarlens}

\subsection*{Your turn}
\begin{itemize}
\raggedright
  \item \emph{Say it:} \emph{cuit / cuite}
  \item \emph{Say it:} \emph{cuit / cuite}, once more
  \item \emph{Say it:} say \emph{cuit / cuite}
  \item \emph{Recall:} say the French for useful, then the French for useless, then say \emph{cuit / cuite} again
\end{itemize}

\begin{tcolorbox}[breakable,colback=teal!4,colframe=teal!35!black,title={Before you move on}]
What does cuit / cuite mean? (\textquotedblleft{}Cooked\textquotedblright{}.) Say it once more.
\end{tcolorbox}
\section[mûr / mûre]{mûr / mûre — ripe}

\label{lesson:FR-C192-mur-mure}



\begin{quote}
Before the new one: say the French for useless, then the French for cooked.
\end{quote}

\subsection*{You'll want to know: mûr / mûre}
\textbf{mûr / mûre} — \textquotedblleft{}ripe\textquotedblright{}.

\begin{grammarlens}[title={What you've built}]
One more word for this chapter.
\end{grammarlens}

\subsection*{Your turn}
\begin{itemize}
\raggedright
  \item \emph{Say it:} \emph{mûr / mûre}
  \item \emph{Say it:} \emph{mûr / mûre}, once more
  \item \emph{Say it:} say \emph{mûr / mûre}
  \item \emph{Recall:} say the French for useless, then the French for cooked, then say \emph{mûr / mûre} again
\end{itemize}

\begin{tcolorbox}[breakable,colback=teal!4,colframe=teal!35!black,title={Before you move on}]
What does mûr / mûre mean? (\textquotedblleft{}Ripe\textquotedblright{}.) Say it once more.
\end{tcolorbox}
\section[sombre]{sombre — dark}

\label{lesson:FR-C192-sombre}



\begin{quote}
Before the new one: say the French for cooked, then the French for ripe.
\end{quote}

\subsection*{You'll want to know: sombre}
\textbf{sombre} — \textquotedblleft{}dark\textquotedblright{}.

\begin{grammarlens}[title={What you've built}]
One more word for this chapter.
\end{grammarlens}

\subsection*{Your turn}
\begin{itemize}
\raggedright
  \item \emph{Say it:} \emph{sombre}
  \item \emph{Say it:} \emph{sombre}, once more
  \item \emph{Say it:} say \emph{sombre}
  \item \emph{Recall:} say the French for cooked, then the French for ripe, then say \emph{sombre} again
\end{itemize}

\begin{tcolorbox}[breakable,colback=teal!4,colframe=teal!35!black,title={Before you move on}]
What does sombre mean? (\textquotedblleft{}Dark\textquotedblright{}.) Say it once more.
\end{tcolorbox}
\section[clair / claire]{clair / claire — clear, light (of a colour)}

\label{lesson:FR-C192-clair-claire}



\begin{quote}
Before the new one: say the French for ripe, then the French for dark.
\end{quote}

\subsection*{You'll want to know: clair / claire}
\textbf{clair / claire} — \textquotedblleft{}clear, light (of a colour)\textquotedblright{}.

\begin{grammarlens}[title={What you've built}]
One more word for this chapter.
\end{grammarlens}

\subsection*{Your turn}
\begin{itemize}
\raggedright
  \item \emph{Say it:} \emph{clair / claire}
  \item \emph{Say it:} \emph{clair / claire}, once more
  \item \emph{Say it:} say \emph{clair / claire}
  \item \emph{Recall:} say the French for ripe, then the French for dark, then say \emph{clair / claire} again
\end{itemize}

\begin{tcolorbox}[breakable,colback=teal!4,colframe=teal!35!black,title={Before you move on}]
What does clair / claire mean? (\textquotedblleft{}Clear, light (of a colour)\textquotedblright{}.) Say it once more.
\end{tcolorbox}
\section[lointain / lointaine]{lointain / lointaine — distant}

\label{lesson:FR-C192-lointain-lointaine}



\begin{quote}
Before the new one: say the French for dark, then the French for clear.
\end{quote}

\subsection*{You'll want to know: lointain / lointaine}
\textbf{lointain / lointaine} — \textquotedblleft{}distant\textquotedblright{}.

\begin{grammarlens}[title={What you've built}]
One more word for this chapter.
\end{grammarlens}

\subsection*{Your turn}
\begin{itemize}
\raggedright
  \item \emph{Say it:} \emph{lointain / lointaine}
  \item \emph{Say it:} \emph{lointain / lointaine}, once more
  \item \emph{Say it:} say \emph{lointain / lointaine}
  \item \emph{Recall:} say the French for dark, then the French for clear, then say \emph{lointain / lointaine} again
\end{itemize}

\begin{tcolorbox}[breakable,colback=teal!4,colframe=teal!35!black,title={Before you move on}]
What does lointain / lointaine mean? (\textquotedblleft{}Distant\textquotedblright{}.) Say it once more.
\end{tcolorbox}
\section[(dialogue)]{Words from chapters 184-187 — a first pass}

\label{lesson:FR-R192-first-pass-words-from-chapters-184-187}



\begin{quote}
Say the words for clear and distant.
\end{quote}

\subsection*{Your turn}
Say these aloud:

\begin{itemize}
\raggedright
  \item to stir, to pick, to defend, to disappear, to borrow
  \item a muscle, a thumb, a wound, a scar, a cold
  \item friendship, hope, trust, courage, shame
  \item in that case, late, suddenly, until, since
\end{itemize}

\begin{tcolorbox}[breakable,colback=teal!4,colframe=teal!35!black,title={Before you move on}]
To stir? (\textbf{bouger}.) To pick? (\textbf{cueillir}.) To defend? (\textbf{défendre}.) To disappear? (\textbf{disparaître}.) To borrow? (\textbf{emprunter}.) A muscle? (\textbf{le muscle}.) A thumb? (\textbf{le pouce}.) A wound? (\textbf{la blessure}.) A scar? (\textbf{la cicatrice}.) A cold? (\textbf{le rhume}.) Friendship? (\textbf{l'amitié}.) Hope? (\textbf{l'espoir}.) Trust? (\textbf{la confiance}.) Courage? (\textbf{le courage}.) Shame? (\textbf{la honte}.) In that case? (\textbf{alors}.) Late? (\textbf{tard}.) Suddenly? (\textbf{soudain}.) Until? (\textbf{jusqu'à}.) Since? (\textbf{depuis}.)
\end{tcolorbox}
\section[(dialogue)]{Words from chapters 188-192 — a second pass}

\label{lesson:FR-R192-second-pass-words-from-chapters-188-192}



\begin{quote}
Say the words for near to and far from.
\end{quote}

\subsection*{Your turn}
Say these aloud:

\begin{itemize}
\raggedright
  \item near to, far from, at the bottom, on the right, on the left
  \item despite, unfortunately, luckily, certainly, really
  \item cheerful, serious, shy, brave, fearful
  \item simple, complicated, necessary, useful, useless
  \item cooked, ripe, dark, clear, distant
\end{itemize}

\begin{tcolorbox}[breakable,colback=teal!4,colframe=teal!35!black,title={Before you move on}]
Near to? (\textbf{près de}.) Far from? (\textbf{loin de}.) At the bottom? (\textbf{au fond}.) On the right? (\textbf{à droite}.) On the left? (\textbf{à gauche}.) Despite? (\textbf{malgré}.) Unfortunately? (\textbf{malheureusement}.) Luckily? (\textbf{heureusement}.) Certainly? (\textbf{certainement}.) Really? (\textbf{vraiment}.) Cheerful? (\textbf{joyeux / joyeuse}.) Serious? (\textbf{sérieux / sérieuse}.) Shy? (\textbf{timide}.) Brave? (\textbf{courageux / courageuse}.) Fearful? (\textbf{peureux / peureuse}.) Simple? (\textbf{simple}.) Complicated? (\textbf{compliqué / compliquée}.) Necessary? (\textbf{nécessaire}.) Useful? (\textbf{utile}.) Useless? (\textbf{inutile}.) Cooked? (\textbf{cuit / cuite}.) Ripe? (\textbf{mûr / mûre}.) Dark? (\textbf{sombre}.) Clear? (\textbf{clair / claire}.) Distant? (\textbf{lointain / lointaine}.)
\end{tcolorbox}
